\documentclass[11pt]{article}
\usepackage[margin=25mm]{geometry}
\usepackage{fontspec,unicode-math}
\setmainfont{Latin Modern Roman}
\setmathfont{Latin Modern Math}
\usepackage{amsmath,amssymb,mathtools}
\usepackage{xurl,hyperref}
\hypersetup{colorlinks=true,urlcolor=blue,linkcolor=blue}
\setcounter{secnumdepth}{0}
\setlength{\parindent}{0pt}
\setlength{\parskip}{5pt}
\setlength{\emergencystretch}{3em}
\newcommand{\sref}[2]{\hyperlink{source:#1:#2}{[S#2]}}
\newcommand{\cataloguescope}[1]{\paragraph{Source question.}#1}
\newcommand{\floor}[1]{\left\lfloor#1\right\rfloor}
\begin{document}
\section{Density and Geometry Recovery from Nearest-Neighbor Graphs}\label{statement}

\subsection{Definitions and mathematical statement}

\paragraph{Shared notation.}Standard number systems are denoted by $\mathbb N,\mathbb Z,\mathbb Q,\mathbb R,\mathbb C$; index ranges, computational models and specialized notation follow the cited problem source and the definitions below.


A directed $k$-nearest-neighbor graph links each sample to its $k$ closest distinct samples. An unweighted adjacency matrix records only adjacency, without distances. Similarities consist of translation, orthogonal transformation, and common rescaling.
\paragraph{Statement recorded in the supplied source.}
Let $X_1,\ldots,X_n$ be i.i.d. samples from a density $p$ on $\mathbb R^d$, and let $A$ be the adjacency matrix of their unweighted $k$-nearest-neighbor graph. Can $p(X_i)$ be estimated up to constant factors from $A$ alone, and can the sample locations be approximately reconstructed up to translation, rotation, and rescaling when $n$ is large? \sref{statement}{2}
\paragraph{Scope qualification.}
The supplied question does not specify directed versus symmetrized adjacency, density regularity, scaling of $k$, or reconstruction-error norm.
\paragraph{Source status.}
The supplied catalogue labels this record partially solved. Its source-status metadata was checked against the supplied snapshot on 2026-10-11. This import does not independently establish current open status or mathematical correctness. The catalogue literature assessment is dated 2026-08-21.
\subsection{Short English statement}
Recover sample densities and geometric positions using only nearest-neighbor adjacency data.
\subsection{Sources}
\begin{itemize}
\item[\hypertarget{source:statement:1}{S1}] Ulam AI and the UnsolvedMath Contributors, UnsolvedMath: A Curated Collection of Open Mathematics Problems, catalogue entry. Catalogue curation and metadata: CC BY 4.0. The supplied source question is retained; definitions and the English summary are editorial additions.
\url{https://huggingface.co/datasets/ulamai/UnsolvedMath}.
\item[\hypertarget{source:statement:2}{S2}] Learning Theory and Approximation (2013), Oberwolfach Reports, DOI 10.4171/owr/2012/31, pp. 1895--1948
\url{https://doi.org/10.4171/owr/2012/31}.
\item[\hypertarget{source:statement:3}{S3}] von Luxburg et al., density estimation from unweighted kNN graphs
\url{https://proceedings.neurips.cc/paper/2013/hash/eae27d77ca20db309e056e3d2dcd7d69-Abstract.html}.
\item[\hypertarget{source:statement:4}{S4}] point localization and density estimation from ordinal kNN graphs
\url{https://arxiv.org/abs/1504.00722}.
\textit{Underlying Oberwolfach Reports material retains its source-specific licence, including CC BY-SA 4.0 where applicable; see THIRD\_PARTY\_NOTICES.md.}
\end{itemize}
\paragraph{Preparation and review.}
Original question attributed to the cited source and UnsolvedMath contributors. Codex (GPT-6; exact serving revision not exposed) assisted with definitions, English summaries, source inspection and formatting. No new human mathematical authorship or independent proof verification is asserted. Catalogue identity checked against the complete supplied registry and existing repository statements; known flagged repetitions and obvious identical conjectures excluded. Full mathematical-equivalence/literature review has not been performed. Statement preparation and final review followed contribution prompts C and D; automated checks verify structure and identity, not open status.
\label{end:statement}
\end{document}
